%==============================================================
% BAB III TARGET PENGGUNA DAN MANFAAT YANG DIHARAPKAN
%==============================================================

\section{Target Pengguna dan Manfaat yang Diharapkan}

\subsection{Target Pengguna Solusi}

TropiCare melayani empat kelompok pengguna dalam alur layanan kesehatan primer:

\begin{enumerate}
    \item \textbf{Masyarakat Umum dan Kelompok Usia Produktif (15--64 Tahun):} 
    Warga yang mengalami gejala infeksi akut atau memiliki faktor risiko penyakit tidak menular (hipertensi, hiperglikemia, obesitas sentral). Kelompok ini sering menunda pemeriksaan karena jam kerja dan antrean panjang FKTP. Lewat TropiCare, pengguna dapat melakukan skrining mandiri melalui peramban ponsel tanpa mengunduh aplikasi dan tanpa membuat akun, serta menjawab pertanyaan proksi fisik (pola menggigil, nadi istirahat, dan perubahan ukuran pakaian) bila tidak memiliki alat ukur medis di rumah \citep{Kemenkes2024}.
    
    \item \textbf{Keluarga Pendamping Pasien Komorbid, Lansia, dan Anak:} 
    Anggota keluarga atau pengasuh (\textit{caregiver}) yang merawat lansia dengan penyakit kronis maupun anak yang sedang demam atau diare. Sistem memeriksa obat yang diminum pasien untuk mencegah bahaya salah obat bebas (seperti larangan NSAID pada anak suspek demam berdarah dan larangan obat penahan diare pada disentri \citep{Alam2022,Koo2007}), sekaligus mengingatkan syarat puasa 10--12 jam sebelum mengantar lansia cek darah di Puskesmas \citep{PERKENI2024}.
    
    \item \textbf{Masyarakat Wilayah Kepulauan, Pedalaman, dan Daerah 3T:} 
    Warga di daerah endemis penyakit tropis terabaikan, seperti Lembah Napu dan Lindu di Sulawesi Tengah (Skistosomiasis) serta NTT dan Papua (Malaria dan Frambusia). Rute Kuning Khusus mengarahkan warga langsung ke Puskesmas penyelenggara program eliminasi beserta kontak WhatsApp/teleponnya, sehingga warga tidak membuang biaya perjalanan ke faskes yang tidak menyimpan stok obat program.
    
    \item \textbf{Tenaga Kesehatan dan Pengelola FKTP (Puskesmas dan Klinik Pratama):} 
    Dokter, perawat, petugas loket pendaftaran, dan kader Posyandu/Posbindu. Petugas memindai kode QR pasien untuk membuka rangkuman keluhan, riwayat obat, dan saran pemeriksaan laboratorium di peramban tanpa perlu mewawancarai dan mengetik ulang keluhan dari awal.
\end{enumerate}

\subsection{Manfaat yang Diharapkan}

Penggunaan TropiCare memberikan tiga manfaat terukur bagi pasien, fasilitas kesehatan, dan pengelola program kesehatan:

\begin{enumerate}
    \item \textbf{Manfaat bagi Pasien dan Keluarga:}
    Pasien memperoleh panduan arah layanan dalam waktu sekitar 3 menit tanpa biaya; mengenali 11 tanda bahaya darurat dalam $<500\text{ ms}$ untuk segera menghubungi 119 atau menuju IGD; menghindari perdarahan lambung akibat NSAID saat demam berdarah serta \textit{toxic megacolon} akibat Loperamide saat disentri; dan datang ke laboratorium FKTP dalam kondisi sudah berpuasa 10--12 jam.

    \item \textbf{Manfaat bagi Tenaga Kesehatan dan Puskesmas:}
    Puskesmas dapat mengurangi tumpukan antrean sebesar 30\%--40\% karena pasien bergejala ringan non-komplikasi (SKDI 4A) teredukasi untuk menjalani perawatan mandiri (Rute Hijau) di rumah. Bagi pasien yang datang ke Puskesmas, pemindaian kode QR rangkuman mempercepat proses pra-triase di loket dan membantu dokter memilih pemeriksaan penunjang yang relevan.

    \item \textbf{Manfaat bagi Dinas Kesehatan dan Program Nasional:}
    Skrining awal faktor risiko PTM (hipertensi, diabetes, dislipidemia, dan risiko kardiovaskular) membantu menekan komplikasi pembiayaan katastropik BPJS Kesehatan. Di tingkat wilayah, Dinas Kesehatan dan Ditjen P2P Kemenkes RI memperoleh statistik agregat tren keluhan per provinsi tanpa identitas langsung (dengan penyembunyian kelompok berjumlah $<10$ untuk mencegah pelacakan balik individu) sesuai prinsip pelindungan data UU No. 27 Tahun 2022 \citep{UUPDP2022}.
\end{enumerate}
